\section{刑法论 - 刑法的功能}
\begin{enumerate}
  \item 知识体系。
        \begin{enumerate}
          \item 刑法论。功能、解释、原则、效力。
          \item \zd{犯罪论}。如何定罪。理解性特别强。解决 2 个问题。
                \begin{enumerate}
                  \item 犯罪是否成立 \yjt 犯罪构成要件。客观、主观。
                  \item 成立后 3 个问题。
                        \begin{enumerate}
                          \item \zd{犯罪形态}。犯罪在时间上的发展问题。预备、中止、未遂、既遂。
                          \item \zd{共同犯罪}。犯罪在空间的分布问题。教唆、帮助、实行。最难的一个问题。
                          \item 犯罪在数量上的计算问题。一罪、数罪。
                        \end{enumerate}
                \end{enumerate}
          \item 刑罚论。如何量刑。地位降低。死记硬背。
        \end{enumerate}
  \item 研究犯罪（定罪）/刑罚（量刑）。
  \item 作用。最高的 2 个概念。
        \begin{enumerate}
          \item \zx{保护法益}。刑法保护的利益。
          \item \zx{保障人权}。疑罪从无。
        \end{enumerate}
\end{enumerate}
